% Source-derived chapter 08: Questions
% Original source: geometric-local-systems-general-curves-isomonodromy
\chapter{Questions}
\label{geometric-local-systems-general-curves-isomonodromy-chapter-08}
\source{geometric-local-systems-general-curves-isomonodromy}

\begin{remark}[Source section]
\label{geometric-local-systems-general-curves-isomonodromy-chapter-08-source}
\source{geometric-local-systems-general-curves-isomonodromy}
\source{geometric-local-systems-general-curves-isomonodromy}
This chapter reproduces the corresponding section of the retrieved
LaTeX source, including its definitions, statements, examples, and
proofs. It has not yet been translated into Lean.
\notready
\end{remark}

% The source section title was: Questions
\label{section:questions}
\source{geometric-local-systems-general-curves-isomonodromy}

We conclude with some open questions related to our results.
\subsection{Bounds}
	
\begin{question}
\source{geometric-local-systems-general-curves-isomonodromy}
\notready\label{question:bounds}
\source{geometric-local-systems-general-curves-isomonodromy}
Is the bound of $2\sqrt{g+1}$ appearing in \autoref{cor:stable},
\autoref{theorem:very-general-VHS} and 
\autoref{theorem:isomonodromic-deformation-CVHS} sharp? If not, can one
explicitly construct low rank geometric variations of Hodge structure with
infinite monodromy on a general curve or general $n$-pointed curve?
Do there exist counterexamples to the above results if one replaces $2
\sqrt{g+1}$ by a linear function of $g$?
\end{question}
We have no reason to believe the bound is sharp. The Kodaira-Parshin trick (as
used in \autoref{section:counterexample}, for example) is one source of
variations of Hodge structure on $\mathscr{M}_{g,n}$ of rank bounded in terms of
$g, n$, but it is not the only one. For example, we expect that the representations constructed in
in \cite{koberda2016quotients} are cohomologically rigid and hence underlie
integral variations of Hodge structure by \cite[Theorem
1.1]{esnault2018cohomologically} and \cite[Theorem 3]{simpson1992higgs}.
Assuming Simpson's motivicity conjecture (\cite[Conjecture,
p.~9]{simpson1992higgs}) these constructions should be geometric in nature, though this is not clear from the construction. Of course it would be extremely interesting to prove that these representations arise from algebraic geometry.

The representations constructed in \cite{koberda2016quotients} have rank
growing exponentially in $g$ \cite[Corollary 4.5]{koberda2016quotients}.
It is natural, given our results, to ask if one can use their methods to produce representations of smaller rank.

We also raise a related question about bounds on maps to the moduli space of
curves.
\begin{question}
\source{geometric-local-systems-general-curves-isomonodromy}
\notready
	\label{question:}
\source{geometric-local-systems-general-curves-isomonodromy}
	Fix an integer $g\geq 2$.
	What is the smallest integer $h \geq 2$ for which the generic genus $g$ curve,
	i.e., the generic fiber of 
	$\mathscr M_{g,1} \to \mathscr M_g$,
	has a
	non-constant map to
	$\mathscr M_h$?
\end{question}
\begin{remark}
\source{geometric-local-systems-general-curves-isomonodromy}
\notready
	\label{remark:}
\source{geometric-local-systems-general-curves-isomonodromy}
	Since a map $C \to \mathscr{M}_h$ corresponds to a family of
smooth curves of genus $h$ over $C$, by considering the associated family of
Jacobians, it follows from \autoref{corollary:abelian-schemes} that $h \geq
\sqrt{g+1}$.
The Kodaira-Parshin trick 
\autoref{proposition:kodaira-parshin} 
does not a priori apply to construct maps from the generic curve to $\mathscr{M}_h$, because as written it produces disconnected covers.
But one can apply a variant where one takes a cover defined by a characteristic quotient of the fundamental group
to show there is some (fairly large) value of $h$ for which the generic genus
$g$ curve has a non-constant map to $\mathscr M_h$.
See \cite[Theorem
1.4]{mcmullen:from-dynamics-on-surfaces-to-rational-points-on-curves} for more
details.
\end{remark}


\subsection{Non-abelian Hodge loci}
Let $(C, x_1, \cdots, x_n)$ be an $n$-pointed curve, $\mathbb{V}$ a
$\mathbb{Z}$-local system on $C\setminus\{x_1, \cdots, x_n\}$ with
quasi-unipotent local monodromy around the $x_i$, and let $(\mathscr{E}, \nabla)$ be the associated flat vector bundle. We refer to the locus $H_{\mathbb{V}}$ in $\mathscr{T}_{g,n}$ where the corresponding isomonodromic deformation of $(\mathscr{E}, \nabla)$ underlies a polarizable variation of Hodge structure as an \emph{non-abelian Hodge locus}. By analogy to the famous result on algebraicity of Hodge loci of Cattani-Deligne-Kaplan \cite{cattani1995locus}, it is natural to ask:
\begin{question}[Compare to {\cite[Conjecture 12.3]{simpson62hodge}}]
\source{geometric-local-systems-general-curves-isomonodromy}
\notready \label{question:non-abelian-Hodge}
\source{geometric-local-systems-general-curves-isomonodromy}
Let $Z$ be an irreducible component of $H_{\mathbb{V}}$. Is the image of $Z$ in $\mathscr{M}_{g,n}$ algebraic? 
\end{question}
This would follow if all $\mathbb{Z}$-local systems which underlie polarizable variations of Hodge structure arise from geometry, which is perhaps a folk conjecture (and is conjectured explicitly in \cite[Conjecture 12.4]{simpson62hodge}). Just as \cite{cattani1995locus} provides evidence for the Hodge conjecture, a positive answer to \autoref{question:non-abelian-Hodge} would provide evidence for this conjecture.

When we refer to an analytically very general curve we mean in the sense
of \autoref{definition:general}.
A positive answer to \autoref{question:non-abelian-Hodge} would allow us to replace this with the usual algebraic notion of a very general curve in \autoref{theorem:very-general-VHS}. 
It seems plausible that one can make this replacement in \autoref{corollary:geometric-local-systems} without requiring input from \autoref{question:non-abelian-Hodge}, using the main result of \cite{cattani1995locus}. 
\bibliographystyle{alpha}
\bibliography{bibliography-isomonodromy}
